\documentclass[11pt]{article}
\usepackage[margin=1in]{geometry}
\usepackage{amsmath,amssymb,amsthm}
\usepackage[T1]{fontenc}
\usepackage{lmodern}
\usepackage[colorlinks=true,linkcolor=blue,citecolor=blue,urlcolor=blue]{hyperref}
\newtheorem{theorem}{Theorem}
\newtheorem{lemma}{Lemma}
\newtheorem*{sourcetheorem}{Source theorem}
\newcommand{\TV}{\operatorname{TV}}
\newcommand{\Prob}{\mathbb P}
\title{A negative answer to polynomial sample complexity\protect\\
for deletion-trace reconstruction}
\author{Kaiwen Jin\\\small\texttt{yui.ui.siki@gmail.com}}
\date{October 12, 2026}
\hypersetup{pdftitle={A negative answer to polynomial sample complexity for deletion-trace reconstruction},pdfauthor={Kaiwen Jin}}

\begin{document}
\maketitle
\begin{abstract}
Can every binary string be recovered from polynomially many independent
deletion traces at a fixed deletion probability? The September 2026 lower
bound of OpenAI answers this question negatively. We record its precise
consequence for AIM problem 321, including randomized decoders with
unrestricted computation and constants depending on the channel.
The argument combines the cited one-trace indistinguishability theorem with
an elementary product bound and two-point testing. The underlying
indistinguishability construction is due to OpenAI.
\end{abstract}

\section{The model and the source result}

Fix a known deletion probability $q\in(0,1)$ and an integer $n\ge2$.
For $x\in\{0,1\}^n$, the law $D_q(x)$ deletes each bit independently with
probability $q$ and lists the surviving bits in their original order.
The output space is
\[
 \mathcal T_n=\bigcup_{j=0}^n\{0,1\}^j,
\]
where the empty word is included. A trace reveals its length and its bits,
but no original positions. The sample is a tuple of independent full traces.
A decoder knows $n,q$, may be randomized, and outputs a length-$n$ binary word.
It must be the same decoder for every unknown input.

Problem 321 of the AIM catalogue~\cite{AIM321} asks whether, for every fixed
$q$, there are finite $C_q,k_q>0$ giving exact recovery with probability at
least $2/3$ on every input, at every length $n\ge2$, from at most
$\lceil C_qn^{k_q}\rceil$ traces. Computation time is unrestricted.
This note explains why the later result of OpenAI~\cite{OpenAI2026} rules out
that guarantee. It does not introduce a new indistinguishability construction.

For laws $P,Q$ on a finite set $S$, use the convention
\[
 \TV(P,Q)=\frac12\sum_{z\in S}|P(z)-Q(z)|.
\]
The complete trace laws are supported on $\mathcal T_n$, so this agrees with
the source's sum over all finite binary words. Define
\[
 \delta_n(q)=
 \min_{\substack{x,y\in\{0,1\}^n\\x\ne y}}
 \TV(D_q(x),D_q(y)).
\]
The minimum exists because the collection of distinct pairs is finite and
nonempty.

\begin{sourcetheorem}[OpenAI, one-trace conclusion of Theorem 1.2]
For a channel with any fixed $q\in(0,1)$ and any real exponent $a>0$,
\begin{equation}\label{eq:source}
 \lim_{n\to\infty} n^a\delta_n(q)=0.
\end{equation}
\end{sourcetheorem}

The cited statement is equation (1.2) on page 2
of~\cite{OpenAI2026}. Its inputs are deterministic binary words, and its
observations are complete traces. The limit runs through all integer
lengths; Section 8 of that paper proves the bound for complete trace laws and pads the
constructed pairs to length $n$. In particular, $q$ remains fixed as $n$
grows. The proof below uses only~\eqref{eq:source} from that paper.

\section{From indistinguishability to reconstruction}

Write $P^{\otimes m}$ for the joint law of $m$ independent samples from $P$.
For $m=0$ the observation space is a singleton.

\begin{lemma}\label{lem:testing}
Suppose two distinct inputs $x,y$ give observation laws $P,Q$ on the same
finite set. For every common randomized decoder $A$,
\begin{equation}\label{eq:testing}
 \frac{\Prob_x(A=x)+\Prob_y(A=y)}2
 \le \frac{1+\TV(P,Q)}2.
\end{equation}
For each nonnegative integer $m$,
\begin{equation}\label{eq:product}
 \TV(P^{\otimes m},Q^{\otimes m})\le m\TV(P,Q).
\end{equation}
\end{lemma}

\begin{proof}
Let $f(z)$ and $g(z)$ be the decoder's conditional probabilities of reporting
$x$ and $y$ on observation $z$. They satisfy $0\le f,g\le1$ and $f+g\le1$.
Consequently,
\[
 \Prob_x(A=x)+\Prob_y(A=y)
 \le \mathbb E_P f+\mathbb E_Q(1-f)
 =1+\sum_z(P(z)-Q(z))f(z).
\]
Because $P-Q$ has total mass zero, the sum of its positive entries is
$\TV(P,Q)$. The last sum is at most that positive mass, which
proves~\eqref{eq:testing}. This calculation already averages all randomness
used by the decoder.

For $m\ge1$, the signed mass functions satisfy the exact identity
\[
 P^{\otimes m}-Q^{\otimes m}
 =\sum_{j=0}^{m-1}
 P^{\otimes j}\otimes(P-Q)\otimes Q^{\otimes(m-1-j)}.
\]
Each summand has $\ell^1$ norm $\|P-Q\|_1$: the other factors are probability
laws with mass one. The triangle inequality and division by two give
\eqref{eq:product}. For $m=0$ both product laws are the same point mass.
\end{proof}

\begin{theorem}\label{thm:negative}
For every fixed $q\in(0,1)$ and all finite real constants $C,k>0$, there is
an integer $n_0\ge2$ such that the following holds. For every integer
$n\ge n_0$, every integer $0\le m\le\lceil Cn^k\rceil$, and every randomized
decoder $A:\mathcal T_n^m\to\{0,1\}^n$, some input $x\in\{0,1\}^n$ satisfies
\[
 \Prob\bigl[A(Y_1,\ldots,Y_m)=x\bigr]<\frac23,
 \qquad Y_1,\ldots,Y_m\overset{\mathrm{iid}}{\sim}D_q(x).
\]
Thus the polynomial-sample guarantee in AIM problem 321 fails at every fixed
deletion probability in $(0,1)$.
\end{theorem}

Here $n$ is the known input length, $m$ is the number of independent complete
traces, and the probability includes the decoder's own randomness. The
threshold $n_0$ may depend on $q,C,k$; no uniformity in the channel is required.

\begin{proof}
Fix $q,C,k$ as stated. At each $n\ge2$ choose distinct $x_n,y_n$ attaining
$\delta_n(q)$. They can be chosen before the sample budget and the decoder.
Since $n^k\ge1$, every permitted budget obeys
\[
 m\le\lceil Cn^k\rceil\le Cn^k+1\le(C+1)n^k.
\]
Lemma~\ref{lem:testing} therefore gives
\[
 \TV\bigl(D_q(x_n)^{\otimes m},D_q(y_n)^{\otimes m}\bigr)
 \le (C+1)n^k\delta_n(q)=:\varepsilon_n.
\]
Equation~\eqref{eq:source} with $a=k$ implies $\varepsilon_n\to0$ at this
same fixed channel. Choose $n_0\ge2$ so that $\varepsilon_n<1/3$ for every
integer $n\ge n_0$. Applying~\eqref{eq:testing} to the full $m$-trace laws,
\[
 \frac{\Prob_{x_n}(A=x_n)+\Prob_{y_n}(A=y_n)}2
 \le\frac{1+\varepsilon_n}2<\frac23.
\]
At least one of these two recovery probabilities is strictly below $2/3$.
The pair works for every permitted $m$ and every common decoder. This proves
the stated quantifiers and contradicts any polynomial guarantee valid at all
lengths.
\end{proof}

\paragraph{AI-use statement.}
OpenAI Codex assisted with source retrieval, drafting, and checking the testing
argument in this note. The multiscale lower-bound theorem is credited to
OpenAI~\cite{OpenAI2026}. Independent mathematical review is requested.

\begin{thebibliography}{9}
\raggedright
\bibitem{AIM321}
AIM catalogue, \emph{321. Polynomial sample complexity for deletion-trace
reconstruction}, repository version
\texttt{c9929805cece044705abaaebe1cebddac1ba11a6}.
\url{https://github.com/MColbrook/AIM/blob/c9929805cece044705abaaebe1cebddac1ba11a6/problems/321-polynomial-trace-reconstruction.md}.

\bibitem{OpenAI2026}
OpenAI, \emph{Quantitative lower bounds for trace reconstruction},
OpenAI Math Release preprint, September 24, 2026.
Theorem 1.2, p.~2, equation (1.2); proof in Section 8 following the
construction in Sections 2--7. Repository version
\texttt{fd4aeeb2ee4fc729c18d98444fed42fd0529eeeb}.
\url{https://github.com/openai/math/blob/fd4aeeb2ee4fc729c18d98444fed42fd0529eeeb/preprints/quantitative-lower-bounds-for-trace-reconstruction-September-24-2026/paper.pdf}.
\end{thebibliography}
\end{document}
